\section{Source Section II.D: Generalized Hydrodynamic Equations}
\label{sec:hydrodynamic-equations}

\sourceeqrange{Eqs. (2.35)--(2.53)}

\subsection{Purpose and Scope}

This section begins the dynamic part of the companion.  It reconstructs, in
typed form, the balance-law bookkeeping behind the first part of source
Section II.D: number density, momentum, total energy, internal energy, entropy
production, and the local role of the reversible stress tensor and reversible
entropy flux.

The scope is deliberately local.  We derive differential balance identities
and finite algebraic checks.  We do not derive Appendix-B scaling, do not
resolve the printed-B3 branch, do not infer unpublished simulation-code
choices, and do not convert the local entropy-flux identity into a global
boundary statement.  The integrated entropy balance in \eqsrc{2.53} is kept as
a boundary- and wall-transfer issue.

\subsection{Typed Objects and Assumptions}

Let \(\Omega\subset\mathbb{R}^d\) be a fixed smooth container with outward
unit normal \(\nu\).  The source first states the boundary setting before the
hydrodynamic equations: controllable boundary temperatures and vanishing
velocity on the container wall.  In this section we work locally in the bulk
and record boundary requirements only when an integrated statement would need
them.

The fields are assumed smooth enough for the product rules below.  A weak
formulation would require function spaces and trace data; that is outside this
companion pass.

\begin{center}
\begin{tabular}{@{}lll@{}}
\toprule
Symbol & Type & Role in Section II.D \\
\midrule
\(n:\Omega\times I\to\mathbb{R}_{>0}\) & scalar field & number density \\
\(\rho=mn\) & scalar density & mass density, with particle mass \(m\) \\
\(v:\Omega\times I\to\mathbb{R}^d\) & vector field & velocity \\
\(T:\Omega\times I\to\mathbb{R}_{>0}\) & scalar field & local temperature \\
\(\hat e\) & scalar density & gradient-augmented internal energy \\
\(e_T\) & scalar density & total energy without wall surface energy \\
\(\hat S\) & scalar density & gradient-augmented entropy density \\
\(\Pi=(\Pi_{ij})\) & tensor field & reversible stress or pressure tensor \\
\(\sigma=(\sigma_{ij})\) & tensor field & dissipative viscous stress \\
\(\lambda\) & coefficient field & thermal conductivity \\
\(g e_z\) & vector & downward gravity convention in the source \\
\(M=M(n,T)\) & coefficient field & square-gradient free-energy coefficient \\
\(p_1\) & scalar field & diagonal gradient pressure contribution \\
\(J_s^{\mathrm{rev}}\) & vector flux & reversible entropy-flux contribution \\
\(J_{\mathrm{GA}}\) & vector flux & alternative capillary energy-flux object \\
\bottomrule
\end{tabular}
\end{center}

Repeated indices are summed.  We write
\[
  \partial_i=\frac{\partial}{\partial x_i},
  \qquad
  (\grad v)_{ij}=\partial_i v_j,
  \qquad
  A:B=A_{ij}B_{ij}.
\]
The source stress convention from Appendix~\ref{app:reversible-stress-derivation}
is
\[
  f_i^{\mathrm{rev}}=-\partial_j\Pi_{ij},
  \qquad
  \hbox{reversible stress power}=-\Pi_{ij}\partial_i v_j .
\]
Because the reversible tensor used below is symmetric, switching the dummy
order in the stress-power contraction does not change the scalar value.

\subsection{Balance-Law, Closure, and Local/Global Status Ledger}

The formulas in this section play different logical roles.  The following
ledger is a reader-facing provenance guide; it prevents a balance law, a
constitutive law, and a local algebraic identity from being read as the same
kind of statement.
\begin{description}
\item[Balance-law inputs.]  Source Eqs. (2.35), (2.36), and (2.39) state
  the local bulk number, momentum, and total-energy balances.  They are not
  derived from the Newtonian viscous closure.
\item[Constitutive closures.]  The viscous stress \eqsrc{2.37}, thermal
  conductivity branch in \eqsrc{2.39}--\eqsrc{2.45}, and sign assumptions
  \(\eta\ge0\), \(\zeta\ge0\), \(\lambda\ge0\), \(T>0\) are physical model
  inputs.  The algebra below checks their consequences but does not derive the
  closures.
\item[Bookkeeping identities.]  The mass-balance reduction, stress-flux
  product rule, and total-minus-kinetic subtraction are local algebraic steps.
  They lead to the internal-energy balance \eqsrc{2.40} once the source
  balance laws and stress signs are fixed.
\item[Reversible residual choice.]  The pressure tensor and reversible entropy
  flux are chosen so the nondissipative residual in \eqsrc{2.43} vanishes
  through \eqsrc{2.44}.  This is local entropy-production bookkeeping, not a
  proof of a global entropy theorem.
\item[Integrated statements.]  The global entropy relation \eqsrc{2.53} also
  needs boundary heat and wall surface-energy transfer.  Local flux-gauge
  identities are therefore kept separate from integrated energy or entropy
  closure.
\end{description}

\subsection{Number-Density Balance}

The number-density balance \eqsrc{2.35} is a conservation law:
\begin{equation}
  \partial_t n+\grad\cdot(nv)=0 .\sourceeqbadge{2.35}
  \label{eq:iid-number-balance}
\end{equation}
Multiplying by the constant particle mass \(m\) gives the mass balance
\begin{equation}
  \partial_t\rho+\grad\cdot(\rho v)=0 .
  \label{eq:iid-mass-balance}
\end{equation}
Equivalently, after expanding \(\grad\cdot(nv)=v\cdot\grad n+n\grad\cdot v\),
the material-derivative forms are
\begin{equation}
  D_t n+n\grad\cdot v=0,
  \qquad
  D_t\rho+\rho\grad\cdot v=0,
  \qquad
  D_t=\partial_t+v\cdot\grad .
  \label{eq:iid-material-density-balance}
\end{equation}
No constitutive assumption is used here.  Equation
\eqref{eq:iid-mass-balance} is the bookkeeping identity needed when the
kinetic-energy equation is extracted from the momentum equation.

\subsection{Momentum Balance and Stress-Divergence Convention}

Source Eq. (2.36), the momentum balance, can be written in component form as
\begin{equation}
  \partial_t(\rho v_i)+\partial_j(\rho v_i v_j)
  =
  -\partial_j\Pi_{ij}
  +\partial_j\sigma_{ij}
  -\rho g\,\delta_{iz}.\sourceeqbadge{2.36}
  \label{eq:iid-momentum-components}
\end{equation}
This equation has three different logical ingredients:
\begin{itemize}
\item the left-hand side is conservative momentum transport;
\item the reversible force density is \(-\partial_j\Pi_{ij}\);
\item the dissipative and body-force terms are \(+\partial_j\sigma_{ij}\) and
  \(-\rho g\,\delta_{iz}\).
\end{itemize}
The sign of the reversible force density is the same sign recorded in
Appendix~\ref{app:reversible-stress-derivation}.  This is why the reversible
stress contributes as \(-\Pi:\grad v\) in the internal-energy balance.

\subsection{Dissipative Stress as a Constitutive Law}

Source Eq. (2.37) supplies the Newtonian viscous closure
\begin{equation}
  \sigma_{ij}
  =
  \eta(\partial_i v_j+\partial_j v_i)
  +
  \left(\zeta-\frac{2\eta}{d}\right)\delta_{ij}\,\grad\cdot v .\sourceeqbadge{2.37}
  \label{eq:iid-dissipative-stress}
\end{equation}
This is not derived from the conservation laws.  It is a constitutive law with
viscosities \(\eta\) and \(\zeta\).  Under the standard sign hypotheses
\(\eta\ge0\) and \(\zeta\ge0\), the source writes the corresponding viscous
production as a nonnegative quadratic form:
\begin{equation}
  \dot\epsilon_v=\sigma_{ij}\partial_i v_j\ge0 .
  \label{eq:iid-viscous-production}
\end{equation}
The finite verification layer checks a simplified quadratic proxy for this
statement.  That proxy is a compatibility check on the algebraic form, not a
software proof of the second law.

\subsection{Viscous Heat Production}

The source then rewrites \eqref{eq:iid-dissipative-stress} into the heat
production formula \eqsrc{2.41}.  Define
\[
  \theta=\grad\cdot v=\partial_kv_k,
  \qquad
  S_{ij}=\partial_i v_j+\partial_j v_i-\frac{2}{d}\delta_{ij}\theta .
\]
The tensor \(S_{ij}\) is symmetric and traceless.  Therefore
\[
  \sigma_{ij}
  =
  \eta S_{ij}+\zeta\delta_{ij}\theta,
\]
because substituting the definition of \(S_{ij}\) recovers
\eqref{eq:iid-dissipative-stress}.  Split the velocity gradient into
symmetric and antisymmetric parts,
\[
  \partial_i v_j=D_{ij}+\Omega_{ij},
  \quad
  D_{ij}=\frac12(\partial_i v_j+\partial_j v_i),
  \quad
  \Omega_{ij}=\frac12(\partial_i v_j-\partial_j v_i).
\]
Since \(S_{ij}\Omega_{ij}=0\), \(S_{ij}\delta_{ij}=0\), and
\[
  D_{ij}=\frac12 S_{ij}+\frac1d\delta_{ij}\theta,
\]
the contraction becomes
\begin{align}
  \dot\epsilon_v
  &=\sigma_{ij}\partial_i v_j
  =
  (\eta S_{ij}+\zeta\delta_{ij}\theta)(D_{ij}+\Omega_{ij})
  \nonumber\\
  &=
  \frac{\eta}{2}S_{ij}S_{ij}+\zeta\theta^2
  \nonumber\\
  &=
  \sum_{ij}\frac{\eta}{2}
  \left(
    \partial_i v_j+\partial_j v_i-\frac{2}{d}\delta_{ij}\grad\cdot v
  \right)^2
  +\zeta(\grad\cdot v)^2
  \ge0 .\sourceeqbadge{2.41}
  \label{eq:iid-viscous-production-expanded}
\end{align}
The inequality uses the constitutive sign assumptions \(\eta\ge0\) and
\(\zeta\ge0\).  The verification scripts check the two-dimensional component
decomposition; they do not derive Newtonian viscosity from the rest of the
model.

\subsection{Total Energy Density and Type Check}

The total energy density in \eqsrc{2.38} is
\begin{equation}
  e_T=\hat e+\frac{1}{2}\rho |v|^2 .\sourceeqbadge{2.38}
  \label{eq:iid-total-energy-density}
\end{equation}
Both summands are energy densities.  The first term already contains the
square-gradient internal-energy contribution from Section II.B; the second is
kinetic energy per volume.  Wall surface energy is not included in
\eqref{eq:iid-total-energy-density}; it reappears in the later integrated
entropy statement.

\subsection{Total Energy Balance and Flux Terms}

With the source stress convention, the total-energy balance \eqsrc{2.39} is
\begin{equation}
  \partial_t e_T
  =
  -\grad\cdot
  \left[
    e_T v+(\Pi-\sigma)\cdot v
  \right]
  +
  \grad\cdot(\lambda\grad T)
  -
  \rho g v_z .\sourceeqbadge{2.39}
  \label{eq:iid-total-energy-balance}
\end{equation}
Here \((\Pi-\sigma)\cdot v\) is the mechanical energy flux associated with the
stress tensor convention in \eqref{eq:iid-momentum-components}, while
\(\grad\cdot(\lambda\grad T)\) is the heat-conduction contribution in the
source sign convention.  The coefficient \(\lambda\) and the Fourier-type
conductive branch are source-stated model inputs in \eqsrc{2.39}; the companion
checks the subsequent product-rule bookkeeping but does not derive the
conductive law from the other balances.

Equation \eqref{eq:iid-total-energy-balance} is local.  It does not by itself
state a complete global energy balance for a container with wall surface
energy, boundary heat input, or prescribed wall temperature branches.

\subsection{Internal Energy from Total Energy and Kinetic Bookkeeping}

The internal-energy balance \eqsrc{2.40} follows by subtracting the
kinetic-energy balance from \eqref{eq:iid-total-energy-balance}.  The algebra
uses \eqref{eq:iid-mass-balance} and
\eqref{eq:iid-momentum-components}.

First, dot the momentum equation with \(v_i\).  The conservative part satisfies
the exact product-rule identity
\begin{align}
  v_i\left[
    \partial_t(\rho v_i)+\partial_j(\rho v_i v_j)
  \right]
  &=
  \partial_t\!\left(\frac{1}{2}\rho |v|^2\right)
  +
  \grad\cdot\!\left(\frac{1}{2}\rho |v|^2 v\right)
  \\
  &\quad+
  \frac12 |v|^2\left[
    \partial_t\rho+\partial_j(\rho v_j)
  \right].
  \label{eq:iid-kinetic-left-with-residual}
\end{align}
The bracketed final factor is the mass-balance residual.  Therefore, when
\eqref{eq:iid-mass-balance} holds,
\begin{equation}
  v_i\left[
    \partial_t(\rho v_i)+\partial_j(\rho v_i v_j)
  \right]
  =
  \partial_t\!\left(\frac{1}{2}\rho |v|^2\right)
  +
  \grad\cdot\!\left(\frac{1}{2}\rho |v|^2 v\right).
  \label{eq:iid-kinetic-left}
\end{equation}
The finite SymPy and Wolfram checks verify both the residual factorization and
the mass-balance-specialized identity in a representative component model.

Second, use the stress-flux product rule
\begin{equation}
  v_i\partial_j A_{ij}
  =
  \partial_j(A_{ij}v_i)-A_{ij}\partial_j v_i .
  \label{eq:iid-stress-flux-product-rule}
\end{equation}
This single rule explains the signs in the kinetic-energy equation:
\begin{center}
\begin{tabular}{@{}p{0.28\linewidth}p{0.29\linewidth}p{0.29\linewidth}@{}}
\toprule
Momentum force & Flux contribution & Stress-power contribution \\
\midrule
\(-\partial_j\Pi_{ij}\) & \(-\partial_j(\Pi_{ij}v_i)\) & \(+\Pi_{ij}\partial_j v_i\) \\
\(+\partial_j\sigma_{ij}\) & \(+\partial_j(\sigma_{ij}v_i)\) & \(-\sigma_{ij}\partial_j v_i\) \\
\bottomrule
\end{tabular}
\end{center}
For symmetric stresses this contraction is equivalent to
\(A_{ij}\partial_i v_j\) after relabeling dummy indices.  Applying
\eqref{eq:iid-stress-flux-product-rule} to the reversible and dissipative
stress terms gives
\begin{align}
  \partial_t\!\left(\frac{1}{2}\rho |v|^2\right)
  +\grad\cdot\!\left(\frac{1}{2}\rho |v|^2v\right)
  &=
  -\grad\cdot(\Pi\cdot v)+\Pi:\grad v
  \nonumber\\
  &\quad
  +\grad\cdot(\sigma\cdot v)-\sigma:\grad v
  -\rho g v_z .
  \label{eq:iid-kinetic-balance}
\end{align}

The subtraction can be checked term by term:
\begin{center}
\begin{tabular}{@{}lll@{}}
\toprule
Term in total-energy balance & Term in kinetic balance & Remainder \\
\midrule
\(\partial_t e_T\) & \(\partial_t(\rho |v|^2/2)\) & \(\partial_t\hat e\) \\
\(\grad\cdot(e_Tv)\) & \(\grad\cdot(\rho |v|^2v/2)\) & \(\grad\cdot(\hat e v)\) \\
\(\grad\cdot(\Pi\cdot v)\) & \(\grad\cdot(\Pi\cdot v)-\Pi:\grad v\) & \(-\Pi:\grad v\) \\
\(-\grad\cdot(\sigma\cdot v)\) & \(-\grad\cdot(\sigma\cdot v)+\sigma:\grad v\) & \(+\sigma:\grad v\) \\
\(-\rho g v_z\) & \(-\rho g v_z\) & \(0\) \\
\bottomrule
\end{tabular}
\end{center}
Subtract \eqref{eq:iid-kinetic-balance} from
\eqref{eq:iid-total-energy-balance} and use
\eqref{eq:iid-total-energy-density}.  The mechanical flux and gravity terms
cancel, leaving
\begin{equation}
  \partial_t\hat e
  =
  -\grad\cdot(\hat e v)
  -
  \Pi:\grad v
  +
  \sigma:\grad v
  +
  \grad\cdot(\lambda\grad T).\sourceeqbadge{2.40}
  \label{eq:iid-internal-energy-balance}
\end{equation}
This is the companion form of \eqsrc{2.40}, with
\(\dot\epsilon_v=\sigma:\grad v\).  The cancellation is a balance-law
bookkeeping identity; it is not a constitutive derivation.

\subsection{Bulk Entropy Rate from the Section II.B Variation}

Section II.B gave the first variation of the bulk entropy functional,
including the boundary term.  In the dynamic setting, set
\(\delta\hat e=(\partial_t\hat e)\,dt\) and
\(\delta n=(\partial_t n)\,dt\), then divide by \(dt\).  This gives
\begin{equation}
  \frac{dS_b}{dt}
  =
  \int_{\Omega}
  \left(
    \frac{1}{T}\partial_t\hat e
    -
    \frac{\hat\mu}{T}\partial_t n
  \right)dx
  -
  \int_{\partial\Omega}
    \frac{M}{T}(\nu\cdot\grad n)\partial_t n\,da .
  \sourceeqbadge{2.42}
  \label{eq:iid-entropy-rate-from-variation}
\end{equation}
The boundary term is displayed because later integrated entropy statements
depend on wall and boundary data.  It is not silently discarded.

\subsection{Substitution into the Bulk Entropy Rate}

The bulk integrand in \eqref{eq:iid-entropy-rate-from-variation} is
\[
  I=
  \frac{1}{T}\partial_t\hat e
  -
  \frac{\hat\mu}{T}\partial_t n .
\]
Substitute \eqref{eq:iid-number-balance} and
\eqref{eq:iid-internal-energy-balance}:
\begin{align}
  I
  &=
  \frac1T
  \left[
    -\grad\cdot(\hat e v)-\Pi_{ij}\partial_i v_j
    +\dot\epsilon_v+\grad\cdot(\lambda\grad T)
  \right]
  +
  \frac{\hat\mu}{T}\grad\cdot(nv).
  \label{eq:iid-entropy-rate-local-substitution}
\end{align}
Apply the product rules
\begin{align}
  -\frac1T\grad\cdot(\hat e v)
  &=
  -\grad\cdot\left(\frac{\hat e v}{T}\right)
  +\hat e\,v\cdot\grad\left(\frac1T\right),\\
  \frac{\hat\mu}{T}\grad\cdot(nv)
  &=
  \grad\cdot\left(nv\frac{\hat\mu}{T}\right)
  -nv\cdot\grad\left(\frac{\hat\mu}{T}\right),\\
  -\frac1T\Pi_{ij}\partial_i v_j
  &=
  -\partial_j\left(v_i\frac{\Pi_{ij}}{T}\right)
  +v_i\partial_j\left(\frac{\Pi_{ij}}{T}\right),
\end{align}
where the last line is the stress-flux product rule used with the symmetric
reversible tensor: explicitly,
\(-\Pi_{ij}\partial_i v_j=-\Pi_{ij}\partial_j v_i\) after exchanging the
dummy indices and using \(\Pi_{ij}=\Pi_{ji}\).  Keeping the divergence terms
separate, the nondissipative
bulk residual is the coefficient of \(v_i\):
\begin{equation}
  R_i
  =
  \partial_j\left(\frac{\Pi_{ij}}{T}\right)
  +
  \hat e\,\partial_i\left(\frac1T\right)
  -
  n\,\partial_i\left(\frac{\hat\mu}{T}\right).
  \label{eq:iid-reversible-entropy-residual}
\end{equation}
After integrating over the bulk and using the source wall/no-slip branch for
the velocity divergence terms, the source writes the bulk part as
\begin{equation}
  \left(\frac{dS_b}{dt}\right)_{\mathrm{bulk}}
  =
  \int_{\Omega}
  \left\{
    v_iR_i
    +
    \frac1T\left[\grad\cdot(\lambda\grad T)+\dot\epsilon_v\right]
  \right\}dx .
  \sourceeqbadge{2.43}
  \label{eq:iid-entropy-rate-bulk-243}
\end{equation}
This is an integrated statement under the stated boundary branch, not a claim
that boundary heat and wall surface-energy transfer have disappeared from the
later global entropy balance.

\subsection{Reversible-Stress Residual Cancellation}

For the reversible part to make no bulk entropy production for arbitrary
bulk velocity, the residual \(R_i\) in
\eqref{eq:iid-reversible-entropy-residual} must vanish.  Thus the source
requires
\begin{equation}
  \partial_j\left(\frac{\Pi_{ij}}{T}\right)
  =
  -\hat e\,\partial_i\left(\frac1T\right)
  +
  n\,\partial_i\left(\frac{\hat\mu}{T}\right).
  \sourceeqbadge{2.44}
  \label{eq:iid-reversible-stress-condition}
\end{equation}
This is the reversible-stress condition.  It is not a dissipative
constitutive inequality and it is not a proof of the second law.  It selects
the reversible stress so the nondissipative bulk residual in
\eqref{eq:iid-entropy-rate-bulk-243} is zero.

\subsection{Thermal Heat Production}

After \eqref{eq:iid-reversible-stress-condition}, the heat-conduction part is
\((1/T)\grad\cdot(\lambda\grad T)\).  The product rule gives
\begin{equation}
  \grad\cdot\left(\frac{\lambda\grad T}{T}\right)
  =
  \frac1T\grad\cdot(\lambda\grad T)
  +
  \lambda\grad T\cdot\grad\left(\frac1T\right).
\end{equation}
Since \(\grad(1/T)=-\grad T/T^2\),
\begin{equation}
  \frac1T\grad\cdot(\lambda\grad T)
  =
  \grad\cdot\left(\frac{\lambda\grad T}{T}\right)
  +
  \lambda\frac{|\grad T|^2}{T^2}.
  \label{eq:iid-heat-product-rule}
\end{equation}
The source therefore writes the thermal production contribution as
\begin{equation}
  \dot\epsilon_{\theta}
  =
  \lambda\frac{|\grad T|^2}{T}\ge0 .
  \sourceeqbadge{2.45}
  \label{eq:iid-thermal-production}
\end{equation}
The nonnegativity statement uses \(T>0\) and \(\lambda\ge0\).  This is the
source's Fourier heat-flux convention in the limited sense used here: the sign
and coefficient are taken from Onuki's heat-conduction branch, and the displayed
calculation only verifies the entropy-production product rule.

\subsection{Gibbs--Duhem Identity Used for the Reversible Stress}

Before giving the pressure tensor, the source records the identity
\begin{equation}
  d(p/T)
  =
  -e\,d(1/T)+n\,d(\mu/T).
  \sourceeqbadge{2.46}
  \label{eq:iid-gibbs-duhem-246}
\end{equation}
Here \(e\), \(s\), \(p\), and \(\mu\) are the local nongradient
thermodynamic fields.  The named Gibbs--Duhem status is source-anchored to
\eqsrc{2.46}; this is the source-stated local Gibbs--Duhem differential
unpacked below from the two local thermodynamic differentials.  No external
thermodynamics formula is imported without this local derivation.
Starting from
\[
  d\mu=-s\,dT+\frac1n\,dp,
  \qquad
  n\mu=e+p-Tns,
\]
compute
\begin{align}
  -e\,d(1/T)+n\,d(\mu/T)
  &=
  -e\,d(1/T)+\frac nT\,d\mu+n\mu\,d(1/T)
  \nonumber\\
  &=
  \frac nT\,d\mu+(n\mu-e)d(1/T)
  \nonumber\\
  &=
  \frac nT\,d\mu+(p-Tns)d(1/T).
\end{align}
Using \(d(1/T)=-dT/T^2\) and the local differential
\(d\mu=-s\,dT+n^{-1}dp\),
\begin{align}
  -e\,d(1/T)+n\,d(\mu/T)
  &=
  \left(\frac1Tdp-\frac{ns}{T}dT\right)
  -\frac{p}{T^2}dT+\frac{ns}{T}dT
  \nonumber\\
  &=
  \frac1Tdp-\frac{p}{T^2}dT
  =
  d(p/T).
  \label{eq:iid-gibbs-duhem-derivation}
\end{align}
The verification scripts mirror this as a coefficient check in the formal
differentials \(dp\) and \(dT\).

\subsection{Choice of Reversible Stress and Reversible Entropy Flux}

The symmetric tensor satisfying the reversible cancellation condition is the
pressure-tensor form
\begin{equation}
  \Pi_{ij}
  =
  (p+p_1)\delta_{ij}
  +
  M(\partial_i n)(\partial_j n).\sourceeqbadge{2.47}
  \label{eq:iid-pressure-tensor}
\end{equation}
The diagonal part \(p_1\) is given in \eqsrc{2.48}.  Appendix A derives the
equivalent thermodynamic form
\begin{equation}
  p_1=n\hat\mu-\hat e+T\hat S-p,\sourceeqbadge{2.48}
  \label{eq:iid-p1-thermodynamic-form}
\end{equation}
and the next two subsections first check the reversible-stress condition and
then expand this scalar field in terms of \(M\), \(n\), and their derivatives.
Appendix~\ref{app:korteweg-primary-crosswalk} later compares only the capillary
part of this pressure-positive tensor with Korteweg's inspected Eq.~(20).  The
comparison is a conditional local constitutive map, not an entropy, energy,
boundary, historical-influence, or complete-PDE equivalence.

\subsection{Direct Check of the Reversible-Stress Condition}

The pressure tensor \eqref{eq:iid-pressure-tensor} must satisfy the residual
condition \eqref{eq:iid-reversible-stress-condition}.  Set
\[
  q=p+p_1=n\hat\mu-\hat e+T\hat S,
  \qquad
  \Pi_{ij}=q\delta_{ij}+M(\partial_i n)(\partial_j n).
\]
Then
\begin{equation}
  \partial_j\left(\frac{\Pi_{ij}}{T}\right)
  =
  \partial_i\left(\frac{q}{T}\right)
  +
  \partial_j\left[\frac{M}{T}(\partial_i n)(\partial_j n)\right].
  \label{eq:iid-pressure-tensor-condition-lhs}
\end{equation}
The scalar factor obeys
\[
  \frac{q}{T}=\frac{n\hat\mu}{T}-\frac{\hat e}{T}+\hat S,
\]
so its derivative is
\begin{align}
  \partial_i\left(\frac{q}{T}\right)
  &=
  n\partial_i\left(\frac{\hat\mu}{T}\right)
  +\frac{\hat\mu}{T}\partial_i n
  -\frac1T\partial_i\hat e
  -\hat e\partial_i\left(\frac1T\right)
  +\partial_i\hat S .
\end{align}
The pointwise spatial form of the Section II.B variational identity gives
\begin{equation}
  \partial_i\hat S
  =
  \frac1T\partial_i\hat e
  -\frac{\hat\mu}{T}\partial_i n
  -
  \partial_j\left[\frac{M}{T}(\partial_i n)(\partial_j n)\right].
  \label{eq:iid-section-iib-pointwise-identity}
\end{equation}
Substituting \eqref{eq:iid-section-iib-pointwise-identity} into
\eqref{eq:iid-pressure-tensor-condition-lhs} cancels the gradient-coordinate
term and leaves
\begin{equation}
  \partial_j\left(\frac{\Pi_{ij}}{T}\right)
  =
  -\hat e\partial_i\left(\frac1T\right)
  +
  n\partial_i\left(\frac{\hat\mu}{T}\right),
\end{equation}
which is exactly \eqref{eq:iid-reversible-stress-condition}.  This is a local
residual cancellation check.  It is not a global entropy theorem and it does
not replace the independent virtual-displacement derivation in
Appendix~\ref{app:reversible-stress-derivation}.

\subsection{Derivative Expansion of the Diagonal Gradient Pressure
\texorpdfstring{\(p_1\)}{p1}}

The scalar \(p_1\) is a pressure-like energy density.  It is a local scalar
field on \(\Omega\times I\), not an integrated functional.  The expansion below
uses the Section II.B hatted densities and generalized chemical potential:
\begin{align}
  \hat e
  &=
  e+\frac{1}{2}K|\grad n|^2,
  \label{eq:iid-p1-ehat-split}\\
  \hat S
  &=
  ns-\frac{1}{2}C|\grad n|^2,
  \label{eq:iid-p1-shat-split}\\
  \hat\mu
  &=
  \mu
  -
  T\grad\cdot\left(\frac{M}{T}\grad n\right)
  +
  \frac{1}{2}M_n|_T|\grad n|^2.
  \label{eq:iid-p1-muhat}
\end{align}
Here \(M_n|_T=(\partial M/\partial n)_T\) is a thermodynamic partial
derivative at fixed temperature.  This notation is not the same as the spatial
derivative \(\grad M\), because in the non-isothermal branch
\[
  \grad M
  =
  M_n|_T\,\grad n
  +
  M_T|_n\,\grad T .
\]

Insert \eqref{eq:iid-p1-ehat-split}--\eqref{eq:iid-p1-muhat} into
\eqref{eq:iid-p1-thermodynamic-form}.  Group the result into nongradient and
gradient pieces:
\begin{align}
  p_1
  &=\underbrace{(n\mu-e+Tns-p)}_{\text{nongradient thermodynamic part}}
  \\
  &\quad
  -nT\grad\cdot\left(\frac{M}{T}\grad n\right)
  +\frac12 nM_n|_T|\grad n|^2
  -\frac12(K+TC)|\grad n|^2.
  \label{eq:iid-p1-substitution-ledger}
\end{align}
The non-gradient thermodynamic terms
cancel by the local pressure identity
\begin{equation}
  p=n\mu-e+Tns .
  \label{eq:iid-p-local-identity}
\end{equation}
The remaining terms are
\begin{align}
  p_1
  &=
  -nT\grad\cdot\left(\frac{M}{T}\grad n\right)
  +
  \frac{1}{2}\left(nM_n|_T-K-TC\right)|\grad n|^2
  \nonumber\\
  &=
  -nT\grad\cdot\left(\frac{M}{T}\grad n\right)
  +
  \frac{1}{2}\left(nM_n|_T-M\right)|\grad n|^2,
  \label{eq:iid-p1-pre-divergence-expanded}
\end{align}
where the second line uses the coefficient relation \(M=CT+K\) from
Section II.B.

Now expand the divergence term as a spatial product rule:
\begin{equation}
  T\grad\cdot\left(\frac{M}{T}\grad n\right)
  =
  M\grad^2 n
  +
  T\grad\left(\frac{M}{T}\right)\cdot\grad n .
  \label{eq:iid-p1-divergence-product-rule}
\end{equation}
Combining \eqref{eq:iid-p1-pre-divergence-expanded} and
\eqref{eq:iid-p1-divergence-product-rule} gives the derivative-expanded form
\begin{equation}
  p_1
  =
  \frac{1}{2}\left(nM_n|_T-M\right)|\grad n|^2
  -
  Mn\grad^2 n
  -
  Tn\,\grad n\cdot\grad\left(\frac{M}{T}\right).\sourceeqbadge{2.48}
  \label{eq:iid-p1-derivative-expanded}
\end{equation}
This is the companion derivation of the first line of \eqsrc{2.48}.  The final
term is a spatial non-isothermal term.  Replacing \(M_n|_T\) by a spatial
quotient such as \(M_x/n_x\) would incorrectly add the temperature-gradient
piece \(M_T|_n T_x/n_x\) in one dimension.

With \eqref{eq:iid-p1-derivative-expanded}, the pressure tensor is still
\eqref{eq:iid-pressure-tensor}.  The dyadic term
\(M(\partial_i n)(\partial_j n)\) supplies the off-diagonal gradient stress,
while \(p_1\) supplies the isotropic diagonal gradient contribution.  Appendix
B later rescales this tensor; the printed-B3 versus dimensional-source branch
is not chosen or corrected here.

\subsection{Pointwise Assembly of the Entropy Flux}

The local form underlying the integrated variation
\eqref{eq:iid-entropy-rate-from-variation} retains the divergence generated by
the gradient term:
\begin{equation}
  \partial_t\hat S
  =
  \frac1T\partial_t\hat e
  -\frac{\hat\mu}{T}\partial_t n
  -\partial_j\left(\frac{M}{T}\partial_t n\,\partial_j n\right).
  \label{eq:iid-local-entropy-time-variation}
\end{equation}
This is the time-path specialization of the Section II.B first variation
\eqref{eq:iib-entropy-variation-boundary}; it is not a new constitutive law.
The product-rule calculation leading to
\eqref{eq:iid-reversible-stress-condition}, followed by that condition, puts
the first two terms of \eqref{eq:iid-local-entropy-time-variation} in the form
\begin{align}
  \frac1T\partial_t\hat e
  -\frac{\hat\mu}{T}\partial_t n
  &=-\partial_j\left[
    \frac{\hat e-n\hat\mu}{T}v_j
    +
    \frac{v_i\Pi_{ij}}{T}
    -\frac{\lambda}{T}\partial_jT
  \right]
  +\frac{\dot\epsilon_v+\dot\epsilon_\theta}{T}.
  \label{eq:iid-entropy-flux-before-gradient-combination}
\end{align}
The flux inside the total divergence must still be combined with the last
term of \eqref{eq:iid-local-entropy-time-variation}.  Let
\(q=n\hat\mu-\hat e+T\hat S\), so that
\(\Pi_{ij}=q\delta_{ij}+M(\partial_i n)(\partial_j n)\).  Componentwise,
\begin{align}
  &\frac{\hat e-n\hat\mu}{T}v_j
  +\frac{v_i\Pi_{ij}}{T}
  +\frac{M}{T}\partial_tn\,\partial_jn
  \nonumber\\
  &\quad=
  \hat S v_j
  +
  \frac{M}{T}
  \left(\partial_tn+v_i\partial_i n\right)\partial_jn
  \nonumber\\
  &\quad=
  \hat S v_j
  +
  \frac{M}{T}(D_tn)\partial_jn
  \nonumber\\
  &\quad=
  \hat S v_j
  -\frac{M}{T}n(\grad\cdot v)\partial_jn,
  \label{eq:iid-entropy-flux-mass-balance-combination}
\end{align}
where the final equality uses the material form of number conservation in
\eqref{eq:iid-material-density-balance}.  Without that balance law, the
remaining flux residual is
\((M/T)[D_tn+n\grad\cdot v]\partial_jn\), so the cancellation is not an
algebraic identity in unconstrained fields.

With \eqref{eq:iid-pressure-tensor}, the local entropy balance can be written
as
\begin{equation}
  \partial_t\hat S
  =
  -\grad\cdot
  \left[
    \hat S v
    -
    \frac{M}{T}n(\grad\cdot v)\grad n
    -
    \frac{\lambda}{T}\grad T
  \right]
  +
  \frac{\dot\epsilon_v+\dot\epsilon_{\theta}}{T}.\sourceeqbadge{2.49}
  \label{eq:iid-entropy-balance}
\end{equation}

\subsection{Reversible Entropy Flux}

Equation \eqref{eq:iid-entropy-balance} identifies three flux contributions:
ordinary advection of \(\hat S\), conductive entropy transport associated with
\(\lambda\grad T\), and a reversible square-gradient flux.  The reversible
part is the vector field
\begin{equation}
  J_s^{\mathrm{rev}}
  =
  -\frac{M}{T}n(\grad\cdot v)\grad n .\sourceeqbadge{2.49}
  \label{eq:iid-reversible-entropy-flux}
\end{equation}
It has the type of entropy flux: entropy per unit area per unit time.  The
factor \(\grad\cdot v\) makes it a local compression/expansion contribution,
and \(\grad n\) gives its direction in density-gradient space.  This flux is
not an additional boundary condition; it is part of the local entropy-balance
bookkeeping after the reversible stress has been chosen.

\subsection{Local Capillary Energy-Flux Gauge}

The alternative capillary energy-flux object discussed around \eqsrc{2.50} is
\begin{equation}
  J_{\mathrm{GA}}
  =
  M n(\grad\cdot v)\grad n,
  \qquad
  J_s^{\mathrm{rev}}=-\frac{1}{T}J_{\mathrm{GA}} .\sourceeqbadge{2.50}
  \label{eq:iid-local-flux-gauge}
\end{equation}
The relation in \eqref{eq:iid-local-flux-gauge} is a local algebraic
conversion between a reversible entropy-flux contribution and an energy-flux
object.  It is useful for comparing formulations, but it is not a global
balance theorem.  If a domain integral is taken, the divergence theorem turns
flux differences into boundary integrals.  Those boundary terms must be
tracked with the wall heat and surface-energy transfer data.

The integrated-balance requirements can be summarized without adding a new
source claim:
\begin{center}
\begin{tabular}{@{}lll@{}}
\toprule
Local object & What integration adds & Current companion status \\
\midrule
\(J_s^{\mathrm{rev}}=-J_{\mathrm{GA}}/T\) & boundary flux term & local identity only \\
\(\lambda\grad T\) & imposed wall heat branch & boundary data required \\
\(e_s\) and \(\sigma_s\) & surface-energy/entropy rate & wall transfer required \\
\eqsrc{2.53} & global entropy bookkeeping & unresolved unless boundary data are stated \\
\bottomrule
\end{tabular}
\end{center}

\subsection{One-Dimensional Gravity Equilibrium}

The gravity comments following \eqsrc{2.50} are static, one-dimensional
specializations of the already stated balance laws.  Take \(v=0\),
\(\sigma=0\), homogeneous temperature, and fields depending only on \(z\).  The
\(i=z\) component of the momentum balance \eqref{eq:iid-momentum-components}
becomes
\[
  0=-\frac{d\Pi_{zz}}{dz}-\rho g,
\]
and hence
\begin{equation}
  \frac{d\Pi_{zz}}{dz}=-\rho g .\sourceeqbadge{2.51}
  \label{eq:iid-static-gravity-stress}
\end{equation}
This is a restricted equilibrium statement, not a dynamic approximation.

In the same branch, \(T\) is constant, so
\eqref{eq:iid-reversible-stress-condition} gives
\[
  \frac1T\frac{d\Pi_{zz}}{dz}
  =
  \frac{n}{T}\frac{d\hat\mu}{dz}.
\]
Combining this with \eqref{eq:iid-static-gravity-stress} and \(\rho=mn\)
yields
\[
  n\frac{d\hat\mu}{dz}=-mn g,
  \qquad
  \frac{d\hat\mu}{dz}=-mg .
\]
Therefore
\begin{equation}
  \hat\mu+mgz=\hbox{constant}.\sourceeqbadge{2.52}
  \label{eq:iid-static-gravity-chemical-potential}
\end{equation}
The finite verification layer checks only this specialization and integration
constant bookkeeping; it does not derive a new gravitational model.

\subsection{Global Entropy Balance and Boundary Transfer}

The global statement \eqsrc{2.53} starts from the local entropy balance
\eqref{eq:iid-entropy-balance}, but it also uses boundary and wall data.  Write
the local flux as
\begin{equation}
  J_S
  =
  \hat S v
  -\frac{M}{T}n(\grad\cdot v)\grad n
  -\frac{\lambda}{T}\grad T .
\end{equation}
Integrating \eqref{eq:iid-entropy-balance} over \(\Omega\) gives
\begin{equation}
  \frac{dS_b}{dt}
  =
  -\int_{\partial\Omega}\nu\cdot J_S\,da
  +
  \int_{\Omega}\frac{\dot\epsilon_v+\dot\epsilon_\theta}{T}\,dx .
  \label{eq:iid-integrated-bulk-entropy-before-wall}
\end{equation}
The source wall branch sets the wall velocity to zero and uses the wall density
condition from Section II.C even away from equilibrium.  Thus the boundary
term is not discarded; its capillary density piece is combined with the wall
surface contribution.

The heat-conduction part has a fixed sign in the outward-normal convention:
\begin{equation}
  -\nu\cdot\left(-\frac{\lambda}{T}\grad T\right)
  =
  \frac{\nu\cdot\lambda\grad T}{T} .
  \label{eq:iid-conductive-boundary-sign}
\end{equation}
For the wall contribution, define
\[
  S_s=\int_{\partial\Omega}\sigma_s(n)\,da,
  \qquad
  E_s=\int_{\partial\Omega}e_s(n)\,da,
  \qquad
  \dot e_s=\frac{\partial e_s}{\partial n}\partial_t n .
\]
The wall Helmholtz relation \(f_s=e_s-T\sigma_s\), together with the natural
wall condition from Section II.C,
\[
  M\nu\cdot\grad n+\left(\frac{\partial f_s}{\partial n}\right)_T=0,
\]
converts the capillary density boundary term into the surface-energy rate
term in the total entropy balance.  The resulting source statement is
\begin{equation}
  \frac{dS_{\mathrm{tot}}}{dt}
  =
  \int_{\Omega}
    \frac{\dot\epsilon_v+\dot\epsilon_\theta}{T}\,dx
  +
  \int_{\partial\Omega}
    \frac{\nu\cdot\lambda\grad T+\dot e_s}{T}\,da .\sourceeqbadge{2.53}
  \label{eq:iid-global-entropy-structure}
\end{equation}
This formula is recorded as a source-guided global bookkeeping statement.  It
is not obtained from the local flux identity alone.  The boundary integrand
contains heat transfer through the wall and the time rate of wall surface
energy.  Dropping either contribution changes the global statement unless a
separate boundary branch makes that term vanish.

Thus the companion classification has two parts:
\begin{itemize}
\item the local flux gauge is a finite local identity;
\item the global entropy balance is boundary- and wall-transfer dependent.
\end{itemize}
Later Appendix-B and Section III guidance must preserve this separation.

\subsection{Relation to Appendix A}

Appendix~\ref{app:reversible-stress-derivation} derives the stress tensor
\eqref{eq:iid-pressure-tensor} from a virtual-displacement calculation.  The
present section uses that tensor in the dynamic entropy-production
bookkeeping.  These are complementary local calculations:
\begin{itemize}
\item Appendix A identifies the reversible stress from nondissipative virtual
  work and entropy neutrality.
\item Section II.D shows why the same stress and the reversible entropy flux
  remove nondissipative terms from the entropy-production balance.
\end{itemize}
Neither calculation resolves the Appendix-B B3 scaling branch, the
unpublished simulation-code branch, or the global wall/surface transfer map.

\subsection{Verification Checklist}

The companion verification layer for this section checks finite algebraic
statements only:
\begin{enumerate}
\item the kinetic-energy identity obtained by dotting momentum with velocity
  under the mass balance;
\item internal-energy bookkeeping as total energy minus kinetic energy;
\item a simplified nonnegative viscous quadratic proxy under nonnegative
  coefficients;
\item the stress-flux product rule that produces stress power;
\item the local entropy-flux relation \(J_s^{\mathrm{rev}}=-J_{\mathrm{GA}}/T\);
\item the equivalent local conversion \(T J_s^{\mathrm{rev}}+J_{\mathrm{GA}}=0\);
\item an expected nonzero residual when the reversible entropy flux is
  omitted;
\item an expected nonzero residual when the boundary heat and surface-energy
  transfer term is omitted from the global entropy schematic;
\item the derivative expansion of \(p_1\) from its thermodynamic form, including
  the distinction between \(M_n|_T\) and a spatial derivative of \(M\);
\item a regression that global boundary entropy transfer remains unresolved.
\end{enumerate}
These checks do not prove the hydrodynamic PDEs, the second law,
well-posedness, boundary theory, or numerical simulations.

The executable mirrors for the Section II.D calculations are:
\begin{center}
\small
\begin{tabular}{@{}l@{}}
\texttt{scripts/python/section\_iid\_hydrodynamic\_balance\_sympy.py}\\
\texttt{scripts/wolfram/section\_iid\_hydrodynamic\_balance\_checks.wls}\\
\texttt{tests/test\_section\_iid\_hydrodynamic\_balances.py}\\
\texttt{scripts/python/section\_iid\_p1\_pressure\_expansion\_sympy.py}\\
\texttt{scripts/wolfram/section\_iid\_p1\_pressure\_expansion\_checks.wls}\\
\texttt{tests/test\_section\_iid\_p1\_pressure\_expansion.py}
\end{tabular}
\end{center}
The first group checks the kinetic-energy subtraction, stress-flux product
rule, local entropy-flux gauge, and boundary-transfer residuals.  The second
group checks the derivative expansion of \(p_1\), including the distinction
between \(M_n|_T\) and spatial differentiation of \(M(n,T)\).  Both groups
are finite mirrors of the displayed companion bookkeeping.

\subsection{Source-Procedure Completion Status}

The exhaustive-inventory pass assigns the Section II.D derivation-required
rows to local balance-law subtraction, thermodynamic pressure expansion, and
boundary-transfer bookkeeping.  This section now gives explicit companion
coverage for the local rows:
\begin{itemize}
\item Rows (2.35)--(2.40): density, momentum, total-energy, and
  internal-energy balances are expanded with the mass-balance and
  kinetic-energy subtraction ledgers.
\item Rows (2.41)--(2.45): entropy production is treated as
  bookkeeping plus constitutive sign assumptions, not as a software proof of
  the second law.
\item Rows (2.47)--(2.49): the pressure tensor, \(p_1\), and
  reversible entropy flux are connected to Appendix~A, the Eq. (2.44)
  cancellation condition, and the finite symbolic residual checks.
\item Rows (2.50)--(2.52): the local flux-gauge identity and the
  one-dimensional gravity specialization are displayed as restricted local
  calculations.
\item Row (2.53): the integrated entropy statement is separated from the local
  flux-gauge identity and kept boundary/wall-transfer dependent.
\end{itemize}
Thus the local Section II.D source procedures are closed for companion
purposes, while integrated energy/entropy closure remains boundary dependent.

\subsection*{Tagged verification complement}

\OnukiLinkIIDBalances{The tagged Section II.D balance complement} maps the
local balance and entropy-bookkeeping derivation to its finite scripts and
negative controls. \OnukiLinkIIDPressure{The tagged pressure-tensor complement}
separates the \(p_1\) derivative convention from Appendix-B source branches.
Neither link asserts a continuum, global-boundary, or simulation proof.

\subsection{Open Issues Carried Forward}

\begin{itemize}
\item The derivative-expanded \(p_1\) formula has been derived locally in this
  section.  Appendix~\ref{app:scaled-equations-plan} now records the separate
  dimensional-source and printed-B3 scaling branches; this local Section II.D
  calculation does not choose a simulation-code branch.
\item The alternative capillary energy flux and the global boundary terms are
  local/global bookkeeping issues; the boundary contribution is not closed in
  this section.
\item Appendix-B B3 scaling is not used here to derive the Section II.D local
  balances.  It remains a separate branchwise source-printing issue.
\item The unpublished simulation-code branch is not inferred from the printed
  equations.
\end{itemize}
